\subsection{Tecnologías descartadas}
\subsubsection*{Frontend}
\begin{itemize}
    \item \textbf{Angular CLI}: JUSTIFICAR CON ALGUNA REFERENCIA BIBLIOGRÁFICA Y REFERENCIAS AL ANÁLISIS PREVIO DEL CLIENTE. Aunque inicialmente se valoró el uso de Angular como framework frontend, finalmente se optó por React debido a su mayor flexibilidad y adecuación a aplicaciones altamente interactivas y orientadas a la experiencia de usuario. La naturaleza del proyecto (que incluye un componente social con subida de imágenes, gestión dinámica de productos, carrito de compra y un asistente conversacional basado en IA) requería una arquitectura basada en componentes reutilizables y ligeros que facilitara una rápida iteración de interfaz y una mayor libertad en el diseño visual. React ofrece un ecosistema más amplio y modular para la construcción de interfaces modernas, permitiendo seleccionar únicamente las herramientas necesarias sin imponer una estructura rígida, lo que resulta especialmente apropiado para un proyecto con fuerte enfoque en usabilidad y experiencia de usuario.
\end{itemize}